\section{Etapa 3 --- Electr\'onica anal\'ogica y modelo del transductor}
\label{sec:et3}
\periodo{17 de abril --- 20 de mayo de 2026 \quad$\cdot$\quad 9 d\'ias de trabajo registrados}

La etapa m\'as densa en decisiones de dise\~no. Arranca con el modelo del
ge\'ofono y termina con la cadena anal\'ogica que sigue en uso. Contiene los
\textbf{dos pivotes de arquitectura} del proyecto.

\subsection{El problema de partida}

El SM-24 tiene frecuencia natural $f_n = \SI{10}{\hertz}$ y amortiguamiento
$\zeta_0 = 0{,}25$. En la forma en aceleraci\'on su respuesta es

\begin{equation}
H_a(s) = \frac{-G_0\,s}{s^2 + 2\zeta_0\omega_n s + \omega_n^2},
\qquad \omega_n = 2\pi\cdot\SI{10}{\hertz}.
\label{eq:hgeo}
\end{equation}

Por debajo de $f_n$ la respuesta cae a \SI{20}{\decibel} por d\'ecada. El
problema es que la informaci\'on de las capas profundas del suelo vive
justamente ah\'i: mayor profundidad de investigaci\'on significa mayor longitud
de onda, y mayor longitud de onda significa menor frecuencia. Con el
transductor tal cual, la profundidad alcanzable est\'a acotada por el
transductor, no por el suelo ni por la electr\'onica.

\subsection{Primer intento: filtro RC pasivo de entrada (17 de abril)}

El primer ataque fue un filtro RC pasivo antes del amplificador, dise\~nado por
optimizaci\'on num\'erica (\texttt{FiltroPasivoEntrada.m}): se plante\'o la
funci\'on de transferencia conjunta ge\'ofono\,$+$\,RC ---que resulta de tercer
orden--- y se buscaron $R$ y $C$ que aplanaran la respuesta. Se explor\'o en
paralelo una alternativa por \textbf{observador de estado} en Simulink
(\texttt{SM24\_model\_Observer1.slx}).

Ambas cosas se abandonaron el 18 de abril, menos de 28 horas despu\'es.

\descartado{\textbf{Filtro RC pasivo de entrada.} Motivos registrados en la
bit\'acora: \textbf{(1)} la frecuencia de corte queda a merced de la tolerancia
de los componentes ---con \SIrange{5}{10}{\percent} t\'ipico, $f_c$ se corre lo
suficiente como para amplificar de m\'as la resonancia del ge\'ofono en vez de
compensarla---; \textbf{(2)} deriva t\'ermica de los condensadores en trabajo de
campo entre \SI{0}{\celsius} y \SI{40}{\celsius}; \textbf{(3)} \emph{fase no
lineal}, que para MASW es descalificante, porque el m\'etodo mide justamente
velocidad de fase y un error de fase se traduce directo en error de velocidad;
y \textbf{(4)} el PSoC ya trae DSP, as\'i que el filtrado sale gratis en
componentes.}

\pendiente{La bit\'acora del 18/04 presenta estos motivos como
\emph{reconstrucci\'on} a partir del \emph{diff}, no como testimonio directo
---dice literalmente ``las razones probables''--- y en un pasaje habla de un
ge\'ofono de \SI{4.5}{\hertz}, que no es el SM-24. El motivo (3), fase lineal,
es el t\'ecnicamente decisivo y el que sostiene el resto del dise\~no.
\textbf{Confirmar cu\'al fue el motivo real.}}

\subsection{Lo que reemplaz\'o al RC: FIR con decimaci\'on en cascada}

En lugar del RC se dise\~n\'o una cadena de tres etapas FIR de 32 coeficientes
(\texttt{fir1(31,\,0.15)}, ventana de Hamming), cada una seguida de decimaci\'on
$\times 4$ ---decimaci\'on total $\times 64$.

\begin{center}\small
\begin{tabularx}{\textwidth}{@{}l X@{}}
\toprule
\textbf{Propiedad} & \textbf{Por qu\'e importa ac\'a} \\
\midrule
Fase lineal exacta & Todas las frecuencias se retardan lo mismo: no contamina la curva de dispersi\'on. Es la raz\'on principal \\
Estabilidad incondicional & Solo ceros, sin polos \\
Coeficientes exactos & Sin tolerancia, sin deriva t\'ermica, sin envejecimiento \\
Reconfigurable & Cambiar $f_c$ es cambiar una l\'inea, no un componente \\
Costo computacional & Decimar en cascada cuesta $\approx 32(1+\tfrac14+\tfrac1{16})\approx 42$ MAC/muestra equivalentes, contra los \textasciitilde{}2000 coeficientes que exigir\'ia un \'unico filtro de $\times64$ \\
\bottomrule
\end{tabularx}
\end{center}

\seleccionado{\textbf{El principio qued\'o y estructura todo el sistema:} lo que
puede resolverse en digital con fase controlada se resuelve en digital; lo
anal\'ogico se reserva para lo que \emph{solo} puede hacerse antes de
cuantizar ---acondicionar el nivel y extender la banda del transductor.}

\subsection{Segundo intento: compensaci\'on activa}

Con el RC descartado, el problema de la banda baja segu\'ia intacto: ning\'un
filtro digital puede recuperar informaci\'on que el transductor atenu\'o por
debajo del piso de ruido \emph{antes} del ADC. La extensi\'on de banda tiene que
ocurrir en anal\'ogico.

Entre el 22 y el 30 de abril se probaron tres l\'ineas:

\begin{center}\small
\begin{tabularx}{\textwidth}{@{}l X l@{}}
\toprule
\textbf{Script} & \textbf{Enfoque} & \textbf{Destino} \\
\midrule
\texttt{compensadorKai.m} & Compensador anal\'itico de forma cerrada & Explorado \\
\texttt{compensadorPazos.m} & B\'usqueda exhaustiva sobre el espacio de componentes & Explorado \\
\texttt{Compensador.m} & Filtro activo VCVS / Sallen-Key & \textbf{Descartado 30/04} \\
\bottomrule
\end{tabularx}
\end{center}

La forma objetivo qued\'o fijada como

\begin{equation}
H_3(s) = \frac{2(\zeta_1-\zeta_0)\,\omega_n\,s}
              {s^2 + 2\zeta_1\omega_n s + \omega_n^2},
\label{eq:objetivo}
\end{equation}

con $\zeta_0=0{,}25$ ---el del ge\'ofono, que el numerador debe reproducir para
cancelar sus polos--- y $\zeta_1$ el amortiguamiento \emph{nuevo}, mucho mayor,
que fija el denominador de la etapa pasa-banda. En el script de abril se us\'o
$\zeta_1=100$ como valor de trabajo; la caracterizaci\'on medida de julio
(Etapa~9) da $\zeta_1 \approx 937$ para el circuito calibrado.

Esta es exactamente la estrategia de Ma \emph{et al.} (2023): \textbf{no se
cancelan polos ni ceros}, se sube el amortiguamiento manteniendo $\omega_n$.
La red completa resulta $1 - H_3(s)$ ---pasa-todo menos pasa-banda---, y \emph{eso
explica por qu\'e la cadena anal\'ogica del proyecto tiene una etapa pasa-banda
seguida de una suma}: el pasa-banda y la resta \textbf{son} la compensaci\'on.

\begin{figure}[!ht]
\centering
\includegraphics[width=0.86\textwidth]{extension_banda.png}
\caption{Qu\'e hace la compensaci\'on. Al elevar $\zeta_1$, los polos complejos
del transductor se separan en dos reales y el inferior queda en
$\omega_n/(2\zeta_1)$. El corte de \SI{-3}{\decibel} pasa de
\SI{8.0}{\hertz} sin compensar a \SI{0.51}{\hertz} con el $\zeta_1=10$ de Ma
\emph{et al.}\ y a \SI{0.010}{\hertz} con el $\zeta_1\approx 937$ que resulta
del denominador BP medido en la tanda calibrada del 21/07. Calculado a partir
de la Ec.~\eqref{eq:hgeo} y $1-H_3$ con $f_n = \SI{10.2046}{\hertz}$.
\textbf{Esta es la respuesta \emph{de dise\~no}}; la muesca realmente medida
resulta \SI{26.3}{\decibel} menos profunda de lo previsto (Etapa~9).}
\label{fig:extension}
\end{figure}

\subsection{El pivote de topolog\'ia: Sallen-Key $\to$ MFB (30 de abril)}

El 30 de abril se escribi\'o \texttt{ComparacionFiltros.m}, que superpone los
Bode de tres funciones: la MFB propuesta, la Sallen-Key vigente y la respuesta
objetivo de la Ec.~\eqref{eq:objetivo}. Ese gr\'afico fue la herramienta de
decisi\'on.

\begin{figure}[!ht]
\centering
\includegraphics[width=0.9\textwidth]{comparacion_topologias.png}
\caption{Reproducci\'on del \texttt{ComparacionFiltros.m} original, con los
mismos valores de componente. Arriba, las tres respuestas; abajo, la
desviaci\'on contra el objetivo. En \SIrange{0.1}{200}{\hertz} la MFB se aparta
\SI{0.19}{\decibel} como m\'aximo y la Sallen-Key \SI{0.26}{\decibel}: en
magnitud son intercambiables. \textbf{El pivote no se decidi\'o por respuesta.}}
\label{fig:topologias}
\end{figure}

\seleccionado{\textbf{Topolog\'ia MFB (\emph{Multiple Feedback}).} Criterio, en
orden de peso:
\begin{enumerate}[itemsep=0.1em,topsep=0.2em]
\item \textbf{Menor sensibilidad de $Q$ y $\omega_0$ a las tolerancias.} Es lo
decisivo: se usan componentes de tolerancia floja y hay que replicar varios
canales apareados. La Sallen-Key realimenta en positivo y por lo tanto
\emph{amplifica} el error de componentes; la MFB usa solo realimentaci\'on
negativa.
\item Mejor atenuaci\'on en alta frecuencia.
\item La baja impedancia de entrada de la MFB no penaliza, porque hay etapas
de ganancia antes de ella.
\item La inversi\'on de fase se revierte en digital sin costo.
\item El producto ganancia-ancho de banda de los operacionales es enorme
frente a las frecuencias Rayleigh: no es limitante.
\end{enumerate}
Dicho de otro modo: \emph{las desventajas cl\'asicas de la MFB son inocuas en
esta banda, y su ventaja ataca el problema real}.}

\descartado{\textbf{Topolog\'ia Sallen-Key / VCVS.} No por respuesta ---la
Fig.~\ref{fig:topologias} muestra que cumpl\'ia--- sino por sensibilidad a
tolerancias, riesgo de oscilaci\'on con ganancia mayor que uno, e impedancia de
salida real no nula del operacional.}

\begin{figure}[!ht]
\centering
\includegraphics[width=\textwidth]{AnalogPathFinal_trim.png}
\caption{La cadena anal\'ogica que result\'o de esta etapa y sigue vigente:
PGA $\to$ pasa-banda (MFB) $\to$ sumador en corriente $\to$ pasa-bajos, con
referencia ajustable por etapa. El pasa-banda y la suma \emph{son} la red de
compensaci\'on $1-H_3$: no son dos bloques que casualmente est\'an en serie.}
\label{fig:cadena}
\end{figure}

El mismo d\'ia entr\'o la otra decisi\'on de la etapa: la \textbf{suma en
corriente} mediante amplificador de transimpedancia. M\'ultiples se\~nales
entran por resistencias a un nodo de tierra virtual y el TIA convierte la suma
de corrientes en tensi\'on. Es un sumador lineal que no carga las fuentes y que
escala agregando una resistencia por entrada ---esa es la etapa SUM de la
cadena definitiva.

\subsection{Medici\'on: MFB contra Sallen-Key en banco}

La decisi\'on por an\'alisis se contrast\'o despu\'es con datos. Se capturaron
conjuntos con las dos topolog\'ias (\texttt{MFB\_DigitalFilterLPF.mat} y
\texttt{Sallen\_Key.mat}) y se construy\'o \texttt{analisis\_espectral.m} (264
l\'ineas) como banco de comparaci\'on:

\begin{itemize}[itemsep=0.2em]
\item \textbf{PSD} cruda y filtrada, cada medici\'on en gris y la media
$\pm\sigma$ en color, con las bandas de inter\'es resaltadas;
\item \textbf{SNR por sesi\'on}, para identificar capturas at\'ipicas y
descartarlas con criterio en vez de a ojo;
\item \textbf{espectrogramas}, que hacen visible el evento transitorio ---un
auto pasando aparece como una mancha de energ\'ia en \SIrange{2}{20}{\hertz}
durante \SIrange{1}{3}{\second}.
\end{itemize}

\dato{\textbf{Vol\'umen experimental de la etapa:} 18 sesiones de captura el 28
de abril, 41 sesiones entre el 2 y el 4 de mayo. El patr\'on de tiempos del 4 de
mayo alterna micro-sesiones de \SIrange{3}{16}{\second} ---cambiar un
componente, reconectar, verificar--- con sesiones largas de
\SIrange{30}{60}{\minute} de captura estable.}

\subsection{Primeras se\~nales reales y primera salida a exteriores}

El 20 de abril lleg\'o la primera se\~nal real del SM-24 a la PC, v\'ia
\texttt{geophone\_scope.m} (657 l\'ineas), un osciloscopio virtual en MATLAB
sobre puerto serie. El 29 de abril el ge\'ofono sali\'o por primera vez al
exterior ---las capturas \texttt{mesitaEstacionamiento*.mat}, unos 13 segundos
de se\~nal cada una, en el estacionamiento de la Facultad: un sitio con fuentes
de vibraci\'on abundantes y predecibles (autos, personas caminando) a pocos
metros del laboratorio.

\subsection{La cadena de firmware toma forma}

Para el 29 de abril el PSoC ya ten\'ia el \emph{pipeline} completo
ADC\,$\to$\,DFB\,$\to$\,FIR\,$\to$\,UART, con:

\begin{itemize}[itemsep=0.2em]
\item cuatro configuraciones de ADC seleccionables;
\item el \textbf{DFB} (filtro digital en hardware del PSoC) filtrando la salida
del $\Delta\Sigma$ antes del FIR por software;
\item un \textbf{notch FIR de \SI{50}{\hertz}} contra la interferencia de red
(\texttt{notch.m});
\item protocolo binario propio hacia la PC: \texttt{0xAA} muestra de 12 bytes,
\texttt{0xCC} confirmaci\'on de estado, \texttt{0xDD} evento de \emph{debug},
y comandos \texttt{[0xBB][CMD][LEN][payload]} en sentido inverso;
\item \emph{heartbeat} autom\'atico cada \textasciitilde{}\SI{2}{\second} para
detectar cuelgues.
\end{itemize}

\descartado{\textbf{El notch de \SI{50}{\hertz} termin\'o retirado} (decisi\'on
posterior, Etapa~9). La interferencia de red es determinista y de forma
conocida, as\'i que conviene \emph{estimarla y restarla} en post-proceso por
ajuste no lineal de m\'inimos cuadrados: eso no introduce retardo de grupo, es
reversible y ---lo importante--- no borra se\~nal s\'ismica \'util, que en
\SI{50}{\hertz} todav\'ia la hay.}

\subsection{Calibraci\'on manual del \emph{offset} (10 de mayo)}

``Calibraci\'on a manita ya funciona''. El \emph{offset} de continua de la
cadena diferencial se cancel\'o desplazando a mano las referencias de los VDAC,
con un protocolo de configuraci\'on PSoC$\leftrightarrow$MATLAB
(\emph{request} \texttt{0xA5}) para leer y fijar par\'ametros en caliente.

\seleccionado{Este es el germen directo de la auto-calibraci\'on de la
Etapa~6: primero se demostr\'o que ajustando las referencias se
cancelaba el \emph{offset}, y despu\'es se cerr\'o el lazo para que lo hiciera el
firmware solo.}

\subsection{Cierre del compensador: del ideal al circuito real (17--20 de mayo)}

Hasta ac\'a el compensador exist\'ia como funci\'on de transferencia. Faltaba
convertirlo en un circuito con componentes comprables y operacionales reales,
y eso ocup\'o tres d\'ias duros ---el mensaje de \emph{commit} del 18 de mayo
fue \emph{``ayudame dios anal\'ogico quiero sucumbir al diablo digital''}.

\textbf{El operacional real entra al modelo.} Se dej\'o de suponer ganancia
infinita: se model\'o el amplificador con $A_0 = \SI{90}{\decibel}$ y polo
dominante en \SI{8}{\mega\hertz}, y se reverific\'o la respuesta del
compensador con esa no idealidad
(\texttt{graficar\_compensador\_csv.py}). Este es el paso que suele delatar un
dise\~no que ``andaba en el papel''.

\textbf{B\'usqueda de componentes.} Se escribieron cuatro variantes de
optimizaci\'on en paralelo, atacando el mismo problema desde \'angulos
distintos. La que qued\'o plantea el espacio de b\'usqueda como continuo en tres
dimensiones $[\log_{10}R_1, \log_{10}C_1, \log_{10}C_2]$ y luego:

\begin{enumerate}[itemsep=0.15em]
\item determina $R_2$ imponiendo la condici\'on de $\omega_0$ exacta ---o sea,
elimina una dimensi\'on por restricci\'on en vez de por b\'usqueda;
\item resuelve las razones del sumador anal\'iticamente por m\'inimos cuadrados;
\item reci\'en al final hace \emph{snap} a valores comerciales y resuelve lo que
no exista con redes en T.
\end{enumerate}

\seleccionado{El orden importa y es la lecci\'on transferible: \emph{optimizar
en el continuo, restringir por f\'isica, discretizar al final}. Discretizar
antes de optimizar deja el \'optimo fuera de la grilla.}

\textbf{La cadena entra al PSoC (20 de mayo).} Con el proyecto
\texttt{DiferencialToSingleEnded} se implement\'o por primera vez la cadena
anal\'ogica completa dentro del PSoC ---el \texttt{main.c} que arranca todos los
bloques anal\'ogicos en la secuencia correcta. A la \'una de la ma\~nana, con
osciloscopio conectado: \emph{``Datos y m\'as datos, empieza la locura''}.

\dato{Se agreg\'o adem\'as un AMUX para poder \emph{medir el modo com\'un} de la
entrada diferencial, y el protocolo se comprimi\'o a 4 bytes multiplexando tipo
y canal en el segundo byte.}

\subsection{Balance de la etapa}

\begin{center}\small
\begin{tabularx}{\textwidth}{@{}l X@{}}
\toprule
\textbf{Se qued\'o} & Compensaci\'on activa por elevaci\'on de $\zeta$ (no cancelaci\'on polo-cero); topolog\'ia MFB; suma en corriente por TIA; filtrado de fase lineal en digital; DFB en hardware; protocolo binario propio; ajuste de \emph{offset} por VDAC \\
\addlinespace[3pt]
\textbf{Se descart\'o} & Filtro RC pasivo de entrada; observador de estado; topolog\'ia Sallen-Key/VCVS; notch de \SI{50}{\hertz} (m\'as tarde) \\
\addlinespace[3pt]
\textbf{Qued\'o abierto} & El \emph{offset} se cancelaba a mano y hab\'ia que repetirlo por placa y por ganancia: no escalaba a un arreglo de varios nodos \\
\bottomrule
\end{tabularx}
\end{center}
